\documentclass[11pt]{article}
\usepackage[T1]{fontenc}
\usepackage[utf8]{inputenc}
\usepackage{lmodern}
\usepackage[margin=1in]{geometry}
\usepackage{amsmath,amssymb,amsthm}
\usepackage{microtype}
\usepackage[hidelinks]{hyperref}
\usepackage{bookmark}
\hypersetup{pdftitle={Global-unitary entanglement extraction: capacity and the exact PPT fidelity exponent},pdfauthor={Yongxian Zhang},pdfsubject={Written analytic proof, not Lean-formalized},pdfkeywords={entanglement extraction, global unitary, complete PPT, fidelity exponent}}
\usepackage{longtable,booktabs,array}
\setlength{\emergencystretch}{2em}
\providecommand{\tightlist}{\setlength{\itemsep}{0pt}\setlength{\parskip}{0pt}}
\setcounter{secnumdepth}{0}
\title{Global-unitary entanglement extraction:\\capacity and the exact PPT fidelity exponent}
\author{Yongxian Zhang\thanks{School of Computer Science and Engineering, South China University of Technology, Guangzhou, China. Email: \href{mailto:mxymmxym1@gmail.com}{mxymmxym1@gmail.com}. ORCID: \href{https://orcid.org/0009-0000-3864-3536}{0009-0000-3864-3536}.}}
\date{October 2026}
\begin{document}
\maketitle
\begin{abstract}
We determine the faithful entanglement-extraction capacity when arbitrary collective global-unitary preprocessing on an existing bipartite system is followed by local channels, LOCC, or a completely PPT channel. For every fixed $2\le m\le n$ and every input spectrum $p$, including zero eigenvalues, these three capacities equal
$\min\{\log m,[\log(mn)-H(p)]/2\}$.
We also determine the exact completely-PPT target-fidelity exponent at every rate. For $0\le R\le\log m$ it is
$\max_{1\le\alpha\le2}(\alpha-1)[2R-\log(mn)+S_\alpha(p)]/\alpha$;
above $\log m$ it is linear with intercept determined by the optimal logarithmic-negativity rate. The proof combines the Rains fidelity-effect SDP with a spectral-prefix approximation, a random Bell-projection construction using the matrix Bernstein inequality, and exact matching by spectral types and escort distributions. A separate deterministic local Kraus construction proves faithful achievability. The global preprocessing is part of the operation class: this is not a formula for ordinary mixed-state LOCC distillation. The exact exponent is claimed for completely PPT channels, not for LOCC. Exact ancillary checks are supplied; the analytic theorem is not Lean-formalized or externally peer reviewed.
\end{abstract}

\hypertarget{task-and-results}{%
\subsection{1. Task and results}\label{task-and-results}}

Fix integers \(2\le m\le n\), a density matrix \(\rho\) on \(\mathbb C^m\otimes\mathbb C^n\), and its eigenvalue probability vector \(p\), with zeros permitted. Set \[
 A=\log m,\qquad B=\log n,\qquad L=A+B.
 \tag{1}
\] At blocklength \(k\) one first applies an arbitrary global unitary \(U_k\) on the existing \(m^k\times n^k\) system. One then applies a deterministic channel to two output systems of dimension \(K\). No ancillary system is supplied before the global unitary. Local ancillas, discarding, and classical communication are permitted in the subsequent LOCC stage. No postselected branch is renormalized as a deterministic success.

Let \(\Phi_K\) be the normalized rank-one maximally entangled projector. For an operation class \(\mathcal C\), define the \textbf{squared-fidelity convention} \[
 f_k^{\mathcal C}(K)=\sup_{U_k,\Lambda\in\mathcal C}
 \operatorname{Tr}\!\left[\Phi_K\Lambda(U_k\rho^{\otimes k}U_k^*)\right].
 \tag{2}
\] We consider local product channels (LO), LOCC, and \textbf{completely PPT channels}: both \(\Lambda\) and \(\Gamma_{\mathrm{out}}\Lambda\Gamma_{\mathrm{in}}\) must be completely positive and trace preserving. Here the partial transposes are taken on the second party. The latter class is denoted \(\mathrm{cPPT}\) to avoid confusion with channels that merely send PPT states to PPT states. LOCC is contained in cPPT: partial-transpose conjugation replaces a product Kraus operator \(A_j\otimes B_j\) by \(A_j\otimes\overline{B_j}\), preserving complete positivity; trace preservation is invariant under the conjugation. The same holds for limits of such channels. The supremum in the cPPT case is a maximum by the finite-dimensional semidefinite formulation below.

For a rate \(R\ge0\) use \(K_k=\lceil e^{kR}\rceil\), with \(K_k=1\) when \(R=0\). Write \(H(p)=-\sum p_i\log p_i\) and \(S_\alpha(p)=(1-\alpha)^{-1}\log\sum p_i^\alpha\) for \(\alpha>1\), extended by \(S_1=H\).

\textbf{Theorem 1 (capacity and strong converse).} The three capacities coincide: \[
 \boxed{C_{\mathrm{LO}}=C_{\mathrm{LOCC}}=C_{\mathrm{cPPT}}
 =\min\left\{\log m,\frac{\log(mn)-H(p)}2\right\}.}
 \tag{3}
\] Capacity means the supremum of rates for which the fidelity in (2) tends to one. Every \(R<C\) is achievable by a single global unitary followed by local product channels; classical communication is unnecessary. For every \(R>C\), even cPPT fidelity tends to zero exponentially. No assertion that fidelity tends to one at \(R=C\) is needed.

\textbf{Theorem 2 (exact cPPT fidelity exponent).} For every \(R\ge0\), the following limit exists: \[
 \boxed{\begin{aligned}
 \mathcal E(R)&=\lim_{k\to\infty}-\frac1k\log f_k^{\mathrm{cPPT}}(K_k)\\
 &=(R-A)_++\max_{1\le\alpha\le2}
 \frac{\alpha-1}{\alpha}\bigl(S_\alpha(p)-c_R\bigr),\\
 c_R&=L-2\min\{R,A\}.
 \end{aligned}}
 \tag{4}
\] The maximand at \(\alpha=1\) is defined to be zero. For \(0\le R\le A\) this is \[
 \mathcal E(R)=\max_{1\le\alpha\le2}
 \frac{\alpha-1}{\alpha}\bigl(2R-\log(mn)+S_\alpha(p)\bigr).
 \tag{5}
\] It is zero for \(R\le C\) and strictly positive for \(R>C\). The exact exponent in (4) is for cPPT; the same expression is a converse bound, not a claimed matching exponent, for LO or LOCC.

Define the preceding logarithmic-negativity rate {[}Z{]} by \[
 \mathcal N(p)=\min_{1\le\alpha\le2}
 \frac{A+(\alpha-1)B+\log\sum_i p_i^\alpha}{\alpha}.
 \tag{6}
\] An equivalent high-rate consequence is \[
 \mathcal E(R)=R-\mathcal N(p)\quad(R\ge A).
 \tag{7}
\] The equality in (7) is established below without assuming that logarithmic negativity is an achievable LOCC yield.

\textbf{Relation to a recent converse.} Lami's quadratic converse {[}L, Corollary 13 and equations (100)--(102){]} already implies the entropy-deficit upper bound for fixed-state cPPT distillation: choosing the auxiliary state \(\tau=I/N\) makes the quadratic form \(N^{-1}\mathrm{id}\) and its logarithmic potential is \(\tfrac12\operatorname{Tr}(\rho\log(N\rho))=\tfrac12(\log N-H(\rho))\), with the same logarithm convention. We do not claim this converse principle as new. Below we derive the finite-block inequalities directly from the Rains constraints, uniformly over the varying collective unitaries, and match them by local capacity achievability and an exact all-rate cPPT fidelity exponent. The cited fixed-input distillation problem has no free global-unitary preprocessing and is not claimed to be solved here.

These are spectrum-optimized extraction statements. The initial global unitary is an entangling resource, so (3) is \textbf{not} a formula for the ordinary distillable entanglement of the unrotated state. No claim is made about the exact finite-copy APPT purity conjecture, APPT=AS, or an efficient circuit implementation.

\hypertarget{the-fidelity-effect-exact-necessity-and-sufficiency}{%
\subsection{2. The fidelity effect: exact necessity and sufficiency}\label{the-fidelity-effect-exact-necessity-and-sufficiency}}

Fix input dimensions \(2\le a\le b\), total dimension \(N=ab\), and \(K\ge2\). Partial transpose is an involution and self-adjoint for the trace pairing. If \(\Lambda\) is cPPT, let \(M=\Lambda^*(\Phi_K)\). Complete positivity and trace preservation give \[
 0\le M\le I,\qquad -I/K\le M^\Gamma\le I/K.
 \tag{8}
\] Indeed \(\Phi_K^\Gamma=F_K/K\), where \(F_K\) is the swap, with eigenvalues \(\pm1\). The adjoint of \(\Gamma_{\mathrm{out}}\Lambda\Gamma_{\mathrm{in}}\) is positive and unital, so it preserves the two order bounds on \(F_K/K\) and maps it to \(M^\Gamma\).

Conversely every \(M\) satisfying (8) defines the CPTP map \[
 \Lambda_M(X)=\operatorname{Tr}(MX)\Phi_K+
 \operatorname{Tr}((I-M)X)\frac{I-\Phi_K}{K^2-1}.
 \tag{9}
\] Let \(Q_\pm=(I\pm F_K)/2\) be the output symmetric and antisymmetric projections. Direct partial transposition gives \[
 (\Gamma_{\mathrm{out}}\Lambda_M\Gamma_{\mathrm{in}})(X)
 =\operatorname{Tr}\!\left[\frac{I+KM^\Gamma}{K(K+1)}X\right]Q_+
 +\operatorname{Tr}\!\left[\frac{I-KM^\Gamma}{K(K-1)}X\right]Q_-.
 \tag{10}
\] Both coefficient effects are positive by (8), so this map is completely positive. It is trace preserving: \(\operatorname{Tr}Q_\pm=K(K\pm1)/2\), and the two resulting effects sum to \(I\). Equation (9) has target overlap \(\operatorname{Tr}(MX)\). This proves the precise Rains characterization used here: \[
 f_{a,b}^{\mathrm{cPPT}}(\lambda,K)
 =\max_{U,M\text{ satisfying }(8)}\operatorname{Tr}(U\operatorname{diag}(\lambda)U^*M).
 \tag{11}
\] The channel in (9) is not in general LOCC. We will not substitute it for the local achievability construction in Section 6.

\hypertarget{converse-inequalities-and-a-one-shot-spectral-envelope}{%
\subsection{3. Converse inequalities and a one-shot spectral envelope}\label{converse-inequalities-and-a-one-shot-spectral-envelope}}

For Hermitian \(X\) on \(a\times b\), \[
 \|X^\Gamma\|_\infty\le a\|X\|_\infty.
 \tag{12}
\] To see this, test \(X^\Gamma\) on a unit vector \(v\). If \(s_j\) are its Schmidt coefficients, the eigenvalues of \((vv^*)^\Gamma\) are \(s_j^2\) and \(\pm s_i s_j\). Its trace norm is \((\sum s_j)^2\le a\). Trace duality proves (12). Thus (8), together with invariance of the Hilbert--Schmidt norm under partial transpose, implies \[
 \|M\|_\infty\le\gamma:=\min\{1,a/K\},\qquad
 \operatorname{Tr}M^2\le N/K^2.
 \tag{13}
\] For \(1<\alpha\le2\) let \(q=\alpha/(\alpha-1)\ge2\). Then \(\operatorname{Tr}M^q\le\gamma^{q-2}\operatorname{Tr}M^2\). Hölder's inequality in (11) yields \[
 f_{a,b}^{\mathrm{cPPT}}(\lambda,K)
 \le\gamma^{(2-\alpha)/\alpha}
       (N/K^2)^{(\alpha-1)/\alpha}
       \left(\sum_i\lambda_i^\alpha\right)^{1/\alpha}.
 \tag{14}
\] For \(\alpha=1\) the bound is \(f\le\gamma\). In particular the factor \(K^2\), not \(K\), in (13) is responsible for the half entropy deficit in the capacity.

Let \(\lambda_1\ge\cdots\ge\lambda_N\), \(L_r=\sum_{i\le r}\lambda_i\), and \[
 v_r=\min\{1,a/K,\sqrt{N/r}/K\},\qquad T_K=\max_{1\le r\le N}v_rL_r.
 \tag{15}
\] The decreasing eigenvalues of \(M\) are at most \(v_r\), by (13). The singular-value trace inequality and \(\lambda_i\le L_i/i\) give \[
 f_{a,b}^{\mathrm{cPPT}}(\lambda,K)\le\sum_i\lambda_i v_i\le H_N T_K,
 \qquad H_N=\sum_{i=1}^N1/i.
 \tag{16}
\] The following matching estimate will supply exact asymptotic exponents.

\textbf{Lemma 3 (one-shot approximation).} For every input spectrum, every \(2\le a\le b\), and every integer \(K\ge2\), \[
 \frac{T_K}{12\log(8N)}\le f_{a,b}^{\mathrm{cPPT}}(\lambda,K)\le H_NT_K.
 \tag{17}
\] The lower bound may restrict the output eigenbasis of the global unitary to a fixed blockwise Bell basis, with only its eigenvalue assignment varying. It is not an assertion that the subsequent cPPT channel is local.

\hypertarget{projection-construction}{%
\subsubsection{3.1 Projection construction}\label{projection-construction}}

We give the construction from {[}Z{]} explicitly. Put \(J=\lfloor b/a\rfloor\), \(N_0=a^2J\), so \(N/2\le N_0\le N\). Partition the first \(aJ\) second-party coordinates into blocks of size \(a\). On each \(a\times a\) block use the generalized Bell basis \[
 \phi_{juv}=\frac1{\sqrt a}\sum_{t=0}^{a-1}e^{2\pi iut/a}
 |t\rangle\,|ja+(t+v\bmod a)\rangle
 \quad(0\le j<J,\ 0\le u,v<a).
 \tag{18}
\] Orthogonality follows from the geometric sum of phases. If \(P_i=\phi_i\phi_i^*\), \(W_i=P_i^\Gamma\), and \(I_0\) is the projection onto these blocks, then \[
 \|W_i\|_\infty=1/a,\quad W_i^2=I_j/a^2,
 \quad \sum_i W_i=I_0,\quad \sum_iW_i^2=I_0.
 \tag{19}
\] Here \(I_j\) is the product-block identity containing \(P_i\). The first two identities follow by local-unitary conjugation from the swap divided by \(a\); each block has \(a^2\) Bell vectors. The third uses \(I_0^\Gamma=I_0\). This is valid for composite \(a\) and does not require \(a\) to divide \(b\).

For an integer \(1\le r\le N_0/8\), select each \(P_i\) independently with probability \(p_0=4r/N_0\le1/2\). Their sum \(P\) is a projection. Its rank has mean \(4r\) and variance at most \(4r\), whence Chebyshev gives \(\Pr\{\operatorname{rank}P<r\}\le4/(9r)\le4/9\). The centered matrices \(X_i=(\xi_i-p_0)W_i\) have norm at most \(1/a\) and variance sum \(p_0(1-p_0)I_0\). Apply the two-sided form of {[}T, Theorem 1.4{]} on the \(N_0\)-dimensional range: \[
 \Pr\{\|P^\Gamma-p_0I_0\|_\infty\ge t\}
 \le2N_0\exp\left(-\frac{t^2}{2(p_0+t/(3a))}\right).
 \tag{20}
\] Let \(\ell=\log(8N_0)>1\) and \(t=2\sqrt{p_0\ell}+2\ell/a\). For \(x=\sqrt{p_0\ell}\) and \(y=\ell/a\), \[
 t^2-2\ell(p_0+t/(3a))=2x^2+\frac{20}3xy+\frac83y^2\ge0.
 \tag{21}
\] The probability in (20) is at most \(1/4\). Since \(4/9+1/4<1\), some projection simultaneously has rank at least \(r\) and \[
 \|P^\Gamma\|_\infty\le p_0+t
 \le12\log(8N)\max\{1/a,\sqrt{r/N}\}.
 \tag{22}
\] For the last estimate use \(p_0\le\sqrt{p_0}\), \(\sqrt\ell\le\ell\), \(N_0\ge N/2\), and \(6\sqrt2+2<12\).

\hypertarget{feasible-channel-and-spectral-prefix}{%
\subsubsection{3.2 Feasible channel and spectral prefix}\label{feasible-channel-and-spectral-prefix}}

For \(P\) as above, put \(c=\|P^\Gamma\|_\infty>0\), \(t_0=\min\{1,(Kc)^{-1}\}\), and \(M=t_0P\). This effect satisfies (8). Place the largest \(\operatorname{rank}P\) eigenvalues of the input on its range, assigning eigenvectors to the selected Bell vectors. Equation (9) then gives fidelity at least \(t_0L_r\). Set \(C_N=12\log(8N)>1\). From (22), \[
 t_0\ge\min\left\{1,\frac{\min(a,\sqrt{N/r})}{KC_N}\right\}
 \ge C_N^{-1}v_r.
 \tag{23}
\] For \(r>N_0/8\), \(v_rL_r<4/K\). A constant product-output channel is LO and has fidelity \(1/K\), which already dominates \(C_N^{-1}v_rL_r\). This also covers small dimensions with no integer in the first range. Maximizing over \(r\) proves (17). The basis used for the unitary can be completed with product vectors on unused coordinates; no spectrum is changed or system discarded before the subsequent channel.

\hypertarget{tensor-powers-and-the-exact-variational-exponent}{%
\subsection{4. Tensor powers and the exact variational exponent}\label{tensor-powers-and-the-exact-variational-exponent}}

Let \(s\) be the number of positive eigenvalues, and let \(q\) range over probability distributions on that fixed support. Write \(D(q\|p)=\sum q_i\log(q_i/p_i)\). For a type with integer counts \(kq_i\), the multinomial cardinality \(r_q=k!/\prod_i(kq_i)!\) satisfies \[
 (k+1)^{-s}e^{kH(q)}\le r_q\le e^{kH(q)}.
 \tag{24}
\] For completeness the probability of that type under law \(q\) is \(r_q e^{-kH(q)}\). It is at most one and is a mode. To check the latter, the likelihood ratio of any competing counts \(h_i\) equals \(\prod_i[(kq_i)!/h_i!](kq_i)^{h_i-kq_i}\) after cancellation of the powers of \(k\). Each factor is at most one by comparing its factorial factors with \(kq_i\); a competing positive count at a zero coordinate has probability zero. There are at most \((k+1)^s\) types, proving (24).

The mass of type \(q\) under \(p^{\otimes k}\) is at least \((k+1)^{-s}e^{-kD(q\|p)}\). The largest \(r_q\) eigenvalues have at least this mass. Applying (17) in dimensions \(a=m^k,b=n^k\) gives, up to factors whose logarithms are \(o(k)\), \[
 f_k^{\mathrm{cPPT}}(K_k) \ge
 \exp\{-k\,[D(q\|p)+\max\{0,R-A,R-L/2+H(q)/2\}]\}.
 \tag{25}
\] More explicitly one may divide the right side by \(24\log(8(mn)^k)(k+1)^s\): \(e^{kR}\le K_k\le2e^{kR}\), and (24) gives the stated weight bound. Approximate an arbitrary \(q\) by types of denominator \(k\); entropy and relative entropy are continuous on the positive support. Therefore \[
 \limsup_k-\frac1k\log f_k^{\mathrm{cPPT}}(K_k)
 \le \min_q\left[D(q\|p)+\max\{0,R-A,R-L/2+H(q)/2\}\right].
 \tag{26}
\] Put \(d=(R-A)_+\) and \(c=c_R=L-2\min(R,A)\). A direct two-case check (\(R\le A\) or \(R\ge A\)) shows that the objective in (26) is \[
 d+D(q\|p)+\tfrac12(H(q)-c)_+.
 \tag{27}
\] On the other hand applying (14) to tensor powers proves \[
 \liminf_k-\frac1k\log f_k^{\mathrm{cPPT}}(K_k)
 \ge d+\max_{1\le\alpha\le2}J_c(\alpha),\qquad
 J_c(\alpha)=\frac{\alpha-1}{\alpha}(S_\alpha(p)-c).
 \tag{28}
\] For \(R=0\), fidelity is exactly one. Otherwise \(K_k\ge2\) eventually and all preceding inequalities apply.

\hypertarget{matching-without-an-assumed-minimax-exchange}{%
\subsubsection{4.1 Matching without an assumed minimax exchange}\label{matching-without-an-assumed-minimax-exchange}}

Write \(Z_\alpha=\sum_{p_i>0}p_i^\alpha\) and \(q_{\alpha,i}=p_i^\alpha/Z_\alpha\). Differentiation gives \[
 J_c'(\alpha)=\frac{H(q_\alpha)-c}{\alpha^2},\qquad
 \frac{d}{d\alpha}H(q_\alpha)=-\alpha\operatorname{Var}_{q_\alpha}(\log p_i)\le0.
 \tag{29}
\] If \(H(p)\le c\), the maximum in (28) is \(J_c(1)=0\). The choice \(q=p\) makes the nonconstant term in (27) zero.

If \(H(q_2)\ge c\), the maximum is at \(\alpha=2\). Choosing \(q=q_2\) gives \[
 D(q_2\|p)+\tfrac12(H(q_2)-c)=\tfrac12(S_2(p)-c)=J_c(2).
 \tag{30}
\] In the remaining case \(H(q_2)<c<H(p)\), choose \(\alpha_*\in(1,2)\) with \(H(q_{\alpha_*})=c\). This maximizes \(J_c\), and is unique unless the positive eigenvalues are all equal; that exceptional case cannot satisfy the strict intermediate inequalities. If \(\ell_\alpha=\sum_iq_{\alpha,i}\log p_i\), then \[
 H(q_\alpha)=-\alpha\ell_\alpha+\log Z_\alpha,\quad
 D(q_\alpha\|p)=(\alpha-1)\ell_\alpha-\log Z_\alpha.
\] Substituting \(H(q_{\alpha_*})=c\) yields \[
 D(q_{\alpha_*}\|p)=J_c(\alpha_*).
 \tag{31}
\] The boundary cases agree. Thus the lower construction (26) matches the converse (28), proving (4) and existence of the limit. In particular we have the fully proved identity \[
 \min_q\{D(q\|p)+\tfrac12(H(q)-c)_+\}
 =\max_{1\le\alpha\le2}J_c(\alpha)\quad(c\ge0).
 \tag{32}
\] Here \(c_R\ge B-A\ge0\). No unproved exchange of extrema, finite grid, or optimizer is used. For \(R\ge A\), substituting \(c_R=B-A\) into (4) and rearranging proves (7).

\hypertarget{strong-converse-and-the-entropy-threshold}{%
\subsection{5. Strong converse and the entropy threshold}\label{strong-converse-and-the-entropy-threshold}}

For \(R\le A\), (29) shows that \(\mathcal E(R)=0\) precisely when \(L-2R\ge H(p)\). If this inequality fails, \(J_c'(1)>0\), so some \(\alpha>1\) gives a strictly positive exponent. For \(R>A\), \(d>0\) and the maximum in (4) is nonnegative. It follows that every rate above (3) has exponentially vanishing cPPT fidelity, and hence exponentially vanishing LOCC fidelity as well.

This proof of the converse uses the complete-PPT order constraints, not an assumed Shannon formula for mixed-state entanglement. A weaker bound derived only by testing a maximally mixed separable input would control \(\operatorname{Tr}M\) by \(N/K\), which would lose the factor of two and would not prove (3).

\hypertarget{achievability-by-local-extraction-after-the-global-unitary}{%
\subsection{6. Achievability by local extraction after the global unitary}\label{achievability-by-local-extraction-after-the-global-unitary}}

For finite dimensions and \(1\le K\le a\), put \(u=\lfloor a/K\rfloor\), \(v=\lfloor b/K\rfloor\), and \(r=uv\). Identify good local coordinate subspaces with \(\mathbb C^K\otimes\mathbb C^u\) and \(\mathbb C^K\otimes\mathbb C^v\). The vectors \[
 \psi_{jl}=K^{-1/2}\sum_{x=0}^{K-1}|x,j\rangle_A|x,l\rangle_B
 \quad(1\le j\le u,\ 1\le l\le v)
 \tag{33}
\] are orthonormal. A global unitary puts the largest \(r\) eigenvalues on these vectors, and completes an orthonormal basis for the other eigenvalues.

Alice's local channel has Kraus operators \(A_j=\sum_x|x\rangle\langle x,j|\) on the good subspace, together with \(|0\rangle\langle z|\) for each unused coordinate \(z\); Bob uses the analogous channel. Their Kraus adjoint products sum to the respective identities, so both channels are trace preserving. Each projector \(\psi_{jl}\psi_{jl}^*\) is mapped exactly to \(\Phi_K\). Contributions from the remaining input eigenvectors are positive, so \[
 f_{a,b}^{\mathrm{LO}}(\lambda,K)\ge L_{\lfloor a/K\rfloor\lfloor b/K\rfloor}.
 \tag{34}
\] There is no postselection, division by a success probability, or hidden availability of pure ancillas before \(U\). The unused coordinates are treated by deterministic local channels.

Fix \(0<R<C\). Then \(R<A\le B\), and \[
 k^{-1}\log\bigl(\lfloor m^k/K_k\rfloor\lfloor n^k/K_k\rfloor\bigr)
 \longrightarrow L-2R>H(p).
 \tag{35}
\] For a direct proof that the corresponding prefix mass tends to one, sample \(X\) with law \(p\) on its positive support. The random variable \(-\log p_X\) is finite and has mean \(H(p)\). The law of large numbers implies that the words whose self-information is at most \(k(H(p)+\epsilon)\) have total mass tending to one. Each such word has probability at least \(e^{-k(H+\epsilon)}\), so there are at most \(e^{k(H+\epsilon)}\) of them. For fixed sufficiently small \(\epsilon>0\), (35) makes this number no larger than the prefix rank in (34), eventually. Thus (34) tends to one. Rate zero is trivial. With the converse this proves all three capacity equalities in (3), including the case \(C=0\).

\hypertarget{consequences-and-boundaries}{%
\subsection{7. Consequences and boundaries}\label{consequences-and-boundaries}}

\hypertarget{a-strict-distinction-from-negativity-growth}{%
\subsubsection{7.1 A strict distinction from negativity growth}\label{a-strict-distinction-from-negativity-growth}}

For \(m=n=d\), the capacity is \[
 C=\log d-\tfrac12H(p),
 \qquad \mathcal N(p)=\log d-\tfrac12S_2(p).
 \tag{36}
\] Unless the positive eigenvalues are all equal, \(H(p)>S_2(p)\) and therefore \(C<\mathcal N(p)\). The entropy inequality follows from Jensen's inequality for \(-\log\) under the probability weights \(p_i\): \(H(p)\ge-\log\sum_i p_i^2\), with equality exactly for a uniform positive support. Thus the previous negativity rate must not be identified with an achievable faithful extraction rate.

For the qutrit APPT example \(p=(3,1,1,1,1,1,1,1,1)/11\), \[
 C=\tfrac12\left(\log9-\log11+\tfrac3{11}\log3\right),\qquad
 \mathcal N(p)=\tfrac12\log(153/121).
 \tag{37}
\] Both are positive and the first is strictly smaller. The previously Lean-certified one-copy APPT property supplies context, not a premise in this proof.

\hypertarget{rank-deficient-and-flat-spectra}{%
\subsubsection{7.2 Rank-deficient and flat spectra}\label{rank-deficient-and-flat-spectra}}

If \(p\) is uniform on \(s\) positive entries, all its Rényi entropies equal \(\log s\). Equations (3)--(4) reduce to \[
 C=\min\{A,\tfrac12\log(mn/s)\},\qquad \mathcal E(R)=(R-C)_+.
 \tag{38}
\] This includes pure inputs (\(s=1\)). For the fully maximally mixed input (\(s=mn\)), the optimum fidelity is exactly \(1/K\): the input stays PPT under every global unitary, a cPPT channel keeps it PPT, and a PPT state's overlap with \(\Phi_K\) is at most \(1/K\). Product output achieves equality.

\hypertarget{an-interior-ruxe9nyi-branch-really-occurs}{%
\subsubsection{7.3 An interior Rényi branch really occurs}\label{an-interior-ruxe9nyi-branch-really-occurs}}

Take \(m=2,n=8\) and \(p=(w,1,1,1,1,0,\ldots,0)/(w+4)\) with \(w=2^{4/3}\). At target rate \(R=\log2\), \(c_R=\log4\). The escort distribution at \(\alpha=3/2\) is \((1/2,1/8,1/8,1/8,1/8)\), whose entropy is exactly \(\log4\). Since the positive eigenvalues are unequal, (29) proves strict monotonicity of escort entropy and a unique interior optimizer. Its exponent is \[
 \mathcal E(\log2)=\log(4+2^{4/3})-\tfrac83\log2>0.
 \tag{39}
\] The positivity also follows from \(H(p)>H(q_{3/2})\) and (29). It is not legitimate to retain only \(\alpha=1,2\) in the general fidelity exponent.

\hypertarget{what-is-and-is-not-resolved}{%
\subsubsection{7.4 What is and is not resolved}\label{what-is-and-is-not-resolved}}

The fixed-dimension, arbitrary-spectrum capacity and the entire cPPT fidelity exponent are resolved for the stated two-stage operation class. The exact ordinary LOCC strong-converse exponent, the below-capacity optimal error exponent, standard distillation without global preprocessing, finite-copy APPT purity, and APPT=AS are different questions and are not claimed here. In particular the word PPT always means the explicitly defined \textbf{complete} partial-transpose-conjugation condition when applied to channels in the theorem.

The proof supplies a finite one-shot approximation (17), a local capacity construction, and a matching exponent argument, not merely numerical evidence. It does not supply polynomial-size global circuits or a search algorithm for the Bell permutation. The old immutable qutrit and negativity publications are not modified.

\hypertarget{references-and-attribution}{%
\subsection{References and attribution}\label{references-and-attribution}}

{[}R{]} Eric M. Rains, \emph{A semidefinite program for distillable entanglement}, IEEE Transactions on Information Theory \textbf{47} (2001), 2921--2933, arXiv:quant-ph/0008047. The fidelity-effect constraints and isotropic-output construction are credited to this work, not claimed as new.

{[}L{]} Ludovico Lami, \emph{On PPT entanglement distillation}, arXiv:2610.12454v1, submitted 8 October 2026. Corollary 13 and equations (100)--(102) give the quadratic converse whose maximally mixed auxiliary state yields the entropy-deficit expression. Credited for that prior converse; not needed as an external premise for the direct finite-block proof here.

{[}T{]} Joel A. Tropp, \emph{User-friendly tail bounds for sums of random matrices}, Foundations of Computational Mathematics \textbf{12} (2012), 389--434, arXiv:1004.4389v7, Theorem 1.4. This proved concentration theorem is an explicit external input.

{[}Z{]} Yongxian Zhang, \emph{The collective-unitary logarithmic-negativity rate of every bipartite spectrum}, repository preprint, immutable \href{https://github.com/mxym/math/releases/tag/collective-unitary-negativity-rate-preprint-v1}{version 1}, source \texttt{3fac265d142cbc5ce48700abce5cced59a1a4b8f}. Sections 3--4 give the Bell-projection mechanism reused and restated here. This companion is a written analytic result, not a Lean theorem.

{[}FWTD{]} Kun Fang, Xin Wang, Marco Tomamichel and Runyao Duan, \emph{Non-asymptotic entanglement distillation}, IEEE Transactions on Information Theory \textbf{65} (2019), 6454--6465, arXiv:1706.06221v3. Related finite-blocklength PPT optimization; it does not supply the spectrum-optimized exponent asserted here.

{[}LLF{]} Zhiwen Lin, Ke Li and Kun Fang, \emph{Exponential Analysis for Entanglement Distillation}, arXiv:2601.10190v2. Their operational definitions must not be conflated with the global-unitary preprocessing and completely PPT channels in this paper; in particular PPT-state-preserving operations in their Example 1 are defined by state preservation, not by complete positivity of partial-transpose conjugation. Used for scope comparison, not as a proof premise.

{[}K{]} Tulja Varun Kondra, Pedro Barrios Hita, Justus Neumann, Hermann Kampermann and Dagmar Bruß, \emph{Fundamental limitations on entanglement extraction from purity}, arXiv:2605.29197v1, Supplemental Lemma 6. Qualitative finite-copy activation was known and is not claimed as a new contribution here.

The elementary pure-subspace packing and typicality argument are standard information-theoretic methods. The contribution claimed here is the local achievability matching the entropy converse under collective global-unitary preprocessing, and the exact, single-letter, all-spectrum cPPT fidelity exponent using the spectral projection estimate. Literature screening is not an exhaustive novelty or priority certification. Independent mathematical review remains necessary; finite ancillary calculations and successful document compilation do not certify the whole analytic proof.


\section*{Verification scope and disclosures}
The mathematical argument uses the stated external matrix Bernstein theorem, together with finite-dimensional matrix analysis. The ancillary program independently checks full rational Choi and partially transposed Choi matrices, deterministic local Kraus maps, small cyclotomic Bell-projector identities, exact constants, type cardinalities, and escort algebra; it also rejects deliberately corrupted inputs and source programs. These checks do not by themselves certify the unbounded analytic theorem. The present work is not Lean-formalized and has not been externally peer reviewed.

This research was conducted independently and received no external funding. AI assistance was used in research exploration, proof development, ancillary programming, and manuscript preparation. The affiliation identifies the author's institution of study and does not imply institutional commissioning or endorsement.

The complete proof, source, and ancillary checks are available in the repository \url{https://github.com/mxym/math}, directory \texttt{research/global-unitary-ppt-extraction}. The preceding qutrit APPT purity theorem and its immutable Lean release are separate results and do not formally certify this paper.
\end{document}
